\documentclass{ctexart}

\usepackage{anyfontsize}
\usepackage{amsmath}
\usepackage{graphicx,subfigure}
\usepackage{bm}
\usepackage{geometry}
\geometry{top=2cm,left=1.5cm, right=1.5cm, bottom=2cm}
\usepackage[shortlabels]{enumitem}

\begin{document}

\title{\textbf{实验一\ 测定冰的熔化热}}
\author{1900011604 张植竣}
\date{2020年9月26日}

\maketitle

\section{实验步骤和测量数据处理}

\subsection{实验步骤}

\noindent 1. 调平并校准天平，测定室温$\theta\,$;

\noindent 2. 测定内筒+搅拌器的质量$m_1$;

\noindent 3. 测定外筒+绝热架+绝热盖的质量$m_2$;

\noindent 4. 在内筒中装入热水，控制水的初温比室温$\theta$高10$\,\sim\,$15\textcelsius;

\noindent 5. 测定量热器+水的质量$m_3$;

\noindent 6. 每隔$10\,\mathrm{s}$记录一次水的温度，用线性外推法测定水的初温$T_2$;

\noindent 7. 取出冰，记录冰的初温$T_1$，将冰块放入内筒，记录投冰时刻，同时开始搅拌;

\noindent 8. 边搅拌边每隔$5\,\mathrm{s}$记录一次水的温度，直到水温一段时间内保持恒定不变，记录冰熔化后系统的平衡温度$T_3$;

\noindent 9. 测定量热器+水+冰的质量$m_4$;

\noindent 10. 计算水的质量$m_0 = m_3 - m_2 - m_1$以及冰的质量$m = m_4 - m_3$。

\subsection{测量数据列表}

\begin{table}[ht!]
    \begin{center}
        \begin{tabular}{|c|c|c|}
            \hline
            室温 & $\theta$/\textcelsius & 27.0 \\
            \hline
            冰的初始温度 & $T_1$/\textcelsius & $-$24.3 \\
            \hline
            水的初始温度 & $T_2$/\textcelsius & 36.2 \\
            \hline
            冰熔化后系统的平衡温度 & $T_3$/\textcelsius & 17.9 \\
            \hline
            内筒+搅拌器的质量 & $m_1 \mathrm{/g}$ & 153.73 \\
            \hline
            外筒+绝热架+绝热盖的质量 & $m_2 \mathrm{/g}$ & 548.38 \\
            \hline
            量热器+水的质量 & $m_3 \mathrm{/g}$ & 855.66 \\
            \hline
            量热器+水+冰的质量 & $m_4 \mathrm{/g}$ & 884.71 \\
            \hline
            水的质量 & $m_0 \mathrm{/g}$ & 153.55 \\
            \hline
            冰的质量 & $m \mathrm{/g}$ & 29.05 \\
            \hline
        \end{tabular}
        \caption{测量数据}
    \end{center}
\end{table}

\begin{figure}[ht!]
    \centering
    \includegraphics[width=0.75\linewidth]{Figure_1.png}
    \caption{用线性外推法测量水的初始温度}
\end{figure}

\begin{figure}[ht!]
    \centering
    \includegraphics[width=0.75\linewidth]{Figure_2.png}
    \caption{水的温度变化曲线\newline （蓝色十字表示投冰时刻，橙色线表示室温）}
\end{figure}

\subsubsection*{计算冰的熔化热}

通过下式计算冰的熔化热：

\begin{equation}
    L = \frac{1}{m}(m_0c_0 + m_1c_1)(T_2 - T_3) - c_0(T_3 - T_0) - c_2(T_0 - T_1) \tag{$*$}
\end{equation}

其中$c_0 = 4.18\times10^3\,\mathrm{J/(kg\cdot K)}$，$c_1 = 0.389\times10^3\,\mathrm{J/(kg\cdot K)}$，$c_2 = 1.80\times10^3\,\mathrm{J/(kg\cdot K)}$，$T_0 = 0$\,\textcelsius。

计算得：
\[
    L = 3.23\times10^5\ \mathrm{J/kg}
\]

\subsubsection*{从图中可以得到的信息}

\noindent 1. 投入冰后，热水先迅速降温，是因为水温高、冰的有效面积大，冰熔化快；后缓慢降温，是因为随着冰的不断熔化，冰块逐渐变小，水温逐渐降低；最后趋于稳定，说明此时冰已经完全熔化且熔化后冷水与原有的热水、内筒和搅拌棒相互热交换达到热平衡。

\medskip

\noindent 2. 投冰后开始的约20$\,\sim\,$30\,s内因来不及搅拌，水温比初温低但是几乎没有变化，开始搅拌后温度计示数迅速下降，说明搅拌会迅速加快系统内部热交换速度，使水温下降得更快。

\section{分析与讨论}

\subsection*{主要误差来源}

\noindent 1. 冰的初始温度测量值偏小。由熔化热计算公式($*$)得：

\[\frac{\partial L}{\partial T_1} = c_2 > 0\]

由于从拿出冰到将冰投入量热器的水中之间有一定的时间差，而且冰的温度($\,\sim -20$\textcelsius)与室温($\,\sim 27$\textcelsius)之间有约\,50\textcelsius\,的温差，冰的吸热会十分迅速，导致实际的$T_1$高于从装冰的保温瓶中温度计读出的$T_1$，即$T_1$测量值偏小，故实验得到的$L$偏小。

\medskip

\noindent 2. 冰的质量测量值偏大。冰中含有的杂质、气泡以及运输过程中表面融化的一层水会导致冰质量的测量值偏大，而杂质的热容量相比于冰又可以忽略。于是由熔化热计算公式($*$)得计算出的$L$会偏小。

\medskip

\noindent 3. 温度计零点偏差。由于$T_0$、$T_1$、$T_2$、$T_3$中只有$T_0$是定义值，其余均为测量值，且冰的熔化热公式($*$)中都用的是温度差，温度计零点偏差造成的误差可以由下式估计：

\[\mathrm{d}L = \frac{\partial L}{\partial T_0}\mathrm{d}T_0\]

其中：

\[\frac{\partial L}{\partial T_0} = c_0 - c_2 = 2.38\times10^3\,\mathrm{J/(kg\cdot K)}\]

\subsection*{补偿法的应用和意义}

实验测量中通常有多种因素会造成系统误差，但是不同因素导致的误差可能作用相反，可以用其中一个抵消另一个产生的影响。例如本实验中量热器从外界吸热或放热分别会使$L$测量值偏小或偏大。使用补偿法可以让吸热与放热抵消，提高实验结果的准确性。同样的方法也可以用在测量电阻、电动势上，是实验上的一种重要方法。

\subsection*{实测结果通常偏小的原因}

主要是冰的初始温度测量值偏小以及冰的质量测量值偏大，分析见上文“主要误差来源”部分。

\section{收获与感想}

\noindent 1. 本实验需要尽量使实验过程进行得迅速，对实验者的技能熟练度要求较高。从把冰拿出来、回到座位、打开盖子、放冰，再关上盖子的这个过程中，实验者需要克服冰表面光滑难以拿取、运输过程中冰的吸热、冰粘在毛巾上难以取下、打开量热器的绝热盖时水的散热、放冰时需要立刻记住投冰时间以及水可能溅出、关上绝热盖之后立即开始搅拌等困难。还有就是一边看秒表记录数据另一边要不断搅拌内筒中的热水或冰水混合物。实验者需要大量练习以达到熟练操作的水平。

\medskip

\noindent 2. 在这个实验中，冰的质量误差对结果影响最大。为了减小投冰时水溅出内筒外以及水蒸发后在量热器绝缘盖上结成露滴所造成的误差，考虑将实验步骤改良为测整个量热器的质量以代替测内筒、搅拌器的质量。这样可以有效避免测得量热器+水+冰的质量$m_4$偏小而导致的冰质量测定值偏小。

\medskip

\noindent 3. 需要多次实验确定水初温的合适取值，本次实验中水的初温、末温过低，使得系统与外界之间吸热大于放热，导致测得冰的熔化热偏小。

\end{document}